\section{Polynomial constraints and finite-state variables}\label{sec:extensions}

The box rounding laws need not satisfy additional polynomial inequalities.
For constrained problems we first bound their expected squared violation,
then use a global geometric bound to repair their support. This gives
$O(r^{-\alpha})$ for the constrained preordering and
$O((\log r/r)^\alpha)$ for the constrained ordinary module. Finite-state
variables require a different extension: retain their labels exactly and
smooth only the continuous coordinates. The same separator argument then
gives both box rates without a further label-count factor.

\subsection{Global feasible repair}\label{sec:ext-geometry}

Let the continuous bags and tree be as in \cref{def:junction-tree}, with
$w\ge1$. Assign finitely many polynomial inequalities to the bags and put
\begin{equation}\label{eq:ext-K}
 K=\{x\in[-1,1]^n:g_{bj}(x_{B_b})\ge0\text{ for every }b,j\}\ne\varnothing,
 \qquad V(x)=\sum_{b,j}(-g_{bj}(x_{B_b}))_+^2.
\end{equation}
Here $t_+=\max\{t,0\}$. In this subsection and the next two,
$\fstar=\min_K f$. Equalities can be supplied as opposite inequalities.
A generator need be assigned only once; repeated assignments are counted
repeatedly in our bounds. Define
\[
 d=\max\{1,\max_b\deg f_b,\max_{b,j}\deg g_{bj}\},\qquad
 G=\sum_{b,j}\chebnorm{g_{bj}}\Abud(g_{bj}).
\]
An empty constraint family contributes zero. Let $L_f$ be a Euclidean
Lipschitz constant of $f$ on the entire box; one may take $L_f=\Abud$,
since $\abs{T_k'}\le k^2$.

\begin{assumption}[Global error bound]\label{ass:global-error}
There are $H\ge0$ and $0<\alpha\le1$ such that
\begin{equation}\label{eq:ext-error-bound}
 \dist_2(x,K)\le H V(x)^{\alpha/2}\qquad(x\in[-1,1]^n).
\end{equation}
\end{assumption}

The distance in this assumption is to the common feasible intersection in
all original coordinates. The exponent and constant depend on the supplied
generators, as well as on the geometry of $K$.

\begin{lemma}[Repair of a box law]\label{lem:global-repair}
Under \cref{ass:global-error}, if a box-supported probability law $\sigma$
satisfies $\int V\,d\sigma\le U$, then a probability law $\nu$ supported on
$K$ satisfies
\[
 \int f\,d\nu\le\int f\,d\sigma+L_fH U^{\alpha/2}.
\]
Some point of $K$ has objective no greater than this upper bound.
\end{lemma}
\begin{proof}
For $\varepsilon>0$ choose a finite $\varepsilon$-net of the nonempty
compact set $K$. Map each box point to its nearest net point, breaking ties
by a fixed ordering. This is a Borel map $P_\varepsilon$, and
$\norm{x-P_\varepsilon x}_2\le\dist_2(x,K)+\varepsilon$.
Lipschitz continuity, \eqref{eq:ext-error-bound}, and Jensen's inequality
for the concave function $t^{\alpha/2}$ give
\[
 \int f(P_\varepsilon x)\,d\sigma(x)
 \le\int f\,d\sigma+L_fH(\int V\,d\sigma)^{\alpha/2}+L_f\varepsilon.
\]
The pushforward laws are supported on $K$. Compactness gives a weakly
convergent subsequence as $\varepsilon\downarrow0$; continuity of $f$ gives
the claim in the limit. The minimum on $K$ is at most the resulting mean.
\end{proof}

\begin{proposition}[A common convex Slater point]\label{prop:convex-slater}
Suppose every $g_{bj}$ is concave on the box, and one point $x^0$ in the
box satisfies $g_{bj}(x^0)\ge\sigma_0>0$ for every supplied constraint.
Then \cref{ass:global-error} holds with $\alpha=1$ and
$H=2\sqrt n/\sigma_0$. The objective need not be convex.
\end{proposition}
\begin{proof}
For a box point $x$ put $v=\max_{b,j}(-g_{bj}(x))_+$, with $v=0$ for an
empty family. Set $\theta=v/(\sigma_0+v)$ and
$z=(1-\theta)x+\theta x^0$. Concavity yields
$g_{bj}(z)\ge-(1-\theta)v+\theta\sigma_0=0$, so $z\in K$. Therefore
\[
 \dist_2(x,K)\le\norm{z-x}_2\le2\sqrt n\,\theta
 \le(2\sqrt n/\sigma_0)v\le(2\sqrt n/\sigma_0)\sqrt{V(x)}.
\]
\end{proof}

\begin{example}[Local bounds do not give the global exponent]\label{ex:global-local}
Take bags $\{x,y\}$ and $\{y,z\}$ in the three-dimensional box, with
$g_1=y-x^2$ and $g_2=-y$. Each local set has a linear repair bound with
constant one: increase $y$ to $\max\{y,x^2\}$ for the first, or decrease it
to $\min\{y,0\}$ for the second. Both repairs stay in their bag boxes.
Globally $K=\{x=y=0\}\times[-1,1]$. At $(t,0,0)$ its distance is
$\abs t$, whereas $V=t^4$. Thus a global bound with exponent $\alpha$
requires $\alpha\le1/2$. Both generators are concave and each has a local
strictly feasible point, but they have no common Slater point. This
example concerns repair geometry; it does not assert a hierarchy lower
bound for these generators.
\end{example}

\subsection{Constrained preordering}\label{sec:ext-pre}

For $I\subseteq B_b$ write $q_I=\prod_{i\in I}(1-x_i^2)$. The constrained
preordering relaxation has normalized local functionals through degree
$2r$, exact separator equalities through degree $2r$, and
\begin{equation}\label{eq:ext-pre-localizers}
 L_b(q_Ip^2)\ge0,\qquad L_b(g_{bj}q_Ip^2)\ge0
\end{equation}
whenever the displayed product has total degree at most $2r$. Its moment
value is denoted $\rho_r^{\mathrm{c,pre}}$. Products of two different
additional generators are unnecessary. A full local preordering in all
supplied generators satisfies these conditions; a module using only
$g_{bj}$ times squares need not satisfy the second condition for nonempty
$I$.

\begin{theorem}[Constrained preordering rounding]\label{thm:constr-pre}
Assume \cref{ass:global-error} and
$r\ge\max\{2d,d+w\}$. Put $m=\lfloor(r-d)/w\rfloor+1$ and
$\Dm=2m^2+1$. Every feasible family for
\eqref{eq:ext-pre-localizers} admits a probability law $\nu$ on $K$ with
\begin{equation}\label{eq:ext-pre-bound}
 \int f\,d\nu\le\sum_bL_b(f_b)+\frac{3\Abud}{\Dm}
       +L_fH\left(\frac{12G}{\Dm}\right)^{\alpha/2}.
\end{equation}
The last two terms also bound
$0\le\fstar-\rho_r^{\mathrm{c,pre}}$. For fixed data and width this is
$O(r^{-\alpha})$.
\end{theorem}
\begin{proof}
Use $K_{m,B_b}$ from \cref{lem:pre-density}. Its preordering
certificate has degree at most $2v_b(m-1)\le2(r-d)$. The same lemma gives
nonnegative mass-one densities $h_b(u)=L_b(K_{m,B_b}(x,u))$ with exactly
consistent separators, hence a box law $\sigma$ with these marginals.
As in \cref{thm:pre},
$\abs{\int f\,d\sigma-\sum_bL_b(f_b)}\le3\Abud/\Dm$.

Fix $g=g_{bj}$ and an output $u$. The conditional functional
$F_u(p)=L_b(K_{m,B_b}(x,u)p(x))$ is positive on
$(a+cg)^2$: multiplying the kernel certificate by this square costs at
most $2d$ extra degrees. The second localizers in
\eqref{eq:ext-pre-localizers} also give $F_u(g)\ge0$. Write
$h=F_u(1)\ge0$. If $g(u)<0$, put $a=-g(u)$. Then
$F_u(g-g(u))\ge ah$ and conditional Cauchy--Schwarz gives
$F_u(g-g(u))^2\le hF_u((g-g(u))^2)$.
For $h>0$ these imply $a^2h\le F_u((g-g(u))^2)$; for $h=0$ the same
inequality follows from positivity. If $g(u)\ge0$ the left side is zero.
Consequently, without normalizing by $h$,
\begin{equation}\label{eq:ext-conditional-violation}
 (-g(u))_+^2h_b(u)\le L_b(K_{m,B_b}(x,u)(g(x)-g(u))^2).
\end{equation}

Let $\mathcal J_b$ be the kernel's diagonal Chebyshev operator. Integrating
\eqref{eq:ext-conditional-violation} bounds the violation by
$L_b(\mathcal J_b(g^2)-2g\mathcal J_bg+g^2)$.
This polynomial has degree at most $2\deg g\le2d\le r$, so
\cref{lem:cheb-moment} applies to its coefficients. By
\cref{lem:jackson,lem:cheb-product},
\begin{align*}
 \chebnorm{\mathcal J_b(g^2)-2g\mathcal J_bg+g^2}
 &\le\frac3{\Dm}\bigl[\Abud(g^2)+2\chebnorm g\Abud(g)\bigr]\\
 &\le\frac{12\chebnorm g\Abud(g)}{\Dm}.
\end{align*}
Thus the same global law satisfies $\int V\,d\sigma\le12G/\Dm$.
Apply \cref{lem:global-repair}. A point of $K$ supplies a feasible Dirac
family, so the moment value is at most $\fstar$; taking the infimum over
families proves the gap bound. Finally $m=\Theta_w(r)$ for fixed $d,w$.
\end{proof}

\subsection{Constrained ordinary module}\label{sec:ext-mod}

Here the local conditions are only
\begin{equation}\label{eq:ext-mod-localizers}
 L_b(p^2)\ge0,\quad L_b((1-x_i^2)p^2)\ge0,\quad
 L_b(g_{bj}p^2)\ge0,
\end{equation}
for products of degree at most $2r$, together with normalization and exact
separator consistency. Denote the value by $\rho_r^{\mathrm{c,mod}}$.
Let $d_\infty$ bound all coordinate degrees of objectives and constraints.
Use $s,N,\Psi,\bar\Psi,\omega,\varrho,D,\delta,\Lambda,\eta,\Gamma_v$
and $\Delta_w$ from \cref{sec:ordinary}, with this $d_\infty$, and set
\[
 a_0=1-\delta,\qquad c_s=\frac4{sC_s}\le\frac6{s^2},\qquad
 E=\sum_b C_b\Gamma_{v_b}+\Delta_w\sum_b\osc(f_b),
\]
\begin{equation}\label{eq:ext-U}
 U=\frac{c_s\sum_{b,j}\Abud(g_{bj})^2+
 \sum_{b,j}\chebnorm{g_{bj}}^2(1+\Delta_w-a_0^{v_b})}{1+\Delta_w}.
\end{equation}
The stronger degree reserve below supports conditional positivity and
constraint-difference certificates. The latter are necessary: a
pointwise multivariate Lipschitz inequality alone cannot be evaluated
under a truncated functional.

\begin{lemma}[A constraint-difference certificate]\label{lem:ext-difference}
Write $g=\sum_\beta c_\beta T_\beta$ on a bag, and set
$A_i(g)=\sum_\beta\abs{c_\beta}\beta_i^2$, so
$\sum_i A_i(g)=\Abud(g)$. For every fixed $u$ in the bag box,
\begin{equation}\label{eq:ext-difference}
 \Abud(g)\sum_i A_i(g)(x_i-u_i)^2-(g(x)-g(u))^2
 \in\Mod_{\deg g}(B).
\end{equation}
Every summand of its certificate has degree at most $2\deg g$; for
constant $g$ the polynomial is zero.
\end{lemma}
\begin{proof}
Order the bag coordinates and telescope each tensor Chebyshev term. For
$\ell=(\beta,i)$ with $c_\beta\ne0$ and $k=\beta_i>0$, define
\[
 w_\ell=\abs{c_\beta}k^2,\quad t_\ell=x_i-u_i,\quad
 R_\ell(x)=\operatorname{sign}(c_\beta)
 \frac{T_k(x_i)-T_k(u_i)}{k^2(x_i-u_i)}
 \prod_{j<i}T_{\beta_j}(x_j)\prod_{j>i}T_{\beta_j}(u_j).
\]
The quotient is its polynomial divided-difference extension, of degree
$k-1$. The bound $\abs{T_k'}\le k^2$ shows its normalized modulus is at
most one on the interval. All factors of $R_\ell$ are bounded in modulus
by one, and $\deg R_\ell\le\abs\beta-1$.
For each univariate polynomial factor $a$, the nonnegative polynomial
$1-a^2$ has an interval certificate through degree $2\deg a$ by
\cref{lem:interval-sos}. Constant factors contribute a nonnegative
constant. The identity
$1-\prod_j a_j^2=\sum_j(1-a_j^2)\prod_{k<j}a_k^2$
then gives an ordinary-module certificate for $1-R_\ell^2$ through degree
$2(\abs\beta-1)$: all multipliers are squares and each term has at most
one box generator. Telescoping the original tensor terms gives
$g(x)-g(u)=\sum_\ell w_\ell t_\ell R_\ell(x)$ and
$\sum_\ell w_\ell=\Abud(g)$. With $A=\Abud(g)$ the exact identity
\begin{align*}
 A\sum_\ell w_\ell t_\ell^2-
       \Bigl(\sum_\ell w_\ell t_\ell R_\ell\Bigr)^2
 ={}&A\sum_\ell w_\ell t_\ell^2(1-R_\ell^2)\\
 &+\sum_{\ell<k}w_\ell w_k
        (t_\ell R_\ell-t_kR_k)^2
\end{align*}
proves the claim. Multiplying each complement certificate by $t_\ell^2$
raises its degree to at most $2\abs\beta\le2\deg g$; the remaining
squares satisfy the same degree bound.
\end{proof}

\begin{lemma}[Squared displacement of the SOS kernel]\label{lem:ext-displacement}
Let $j_s(x)=\int(x-u)^2\Psi(x,u)\,d\arcs(u)$. It is globally SOS,
has degree at most $D+2$, and
$c_s-j_s\in\Mod_{(D+2)/2}(\{x\})$. For any ordinary-module-positive
$L_b$ of order $r\ge v_bD+2$ and mass one,
\[
 0\le L_b\Bigl(j_s(x_i)\prod_{\ell\in B_b\setminus\{i\}}
 \omega(x_\ell)\Bigr)\le c_s.
\]
\end{lemma}
\begin{proof}
Put $J_s(x)=\int(x-u)^2\Phi_s(x,u)^2\,d\arcs(u)$, so
$j_s=p_{s,N}J_s$. Integrating polynomial squares yields an SOS polynomial:
its coefficient matrix is the integral of positive semidefinite rank-one
matrices. Thus $J_s$ and $j_s$ are SOS, with degrees at most $2s$ and
$D+2$ respectively. To bound $J_s$, let
$F_s(t)=\sum_{j\in\Z}b_je^{\mathrm{i}jt}$ with
$b_j=(1-\abs j/s)_+$. For $x=\cos\theta$, $u=\cos\phi$,
$\Phi_s(x,u)=[F_s(\theta-\phi)+F_s(\theta+\phi)]/2$.
Convexity of the square and reflection of the $\phi$ integral give
\begin{align*}
 J_s(\cos\theta)
 &\le\int_{-\pi}^{\pi}
 [\cos\theta-\cos(\theta-t)]^2F_s(t)^2\,\frac{dt}{2\pi}\\
 &\le\int_{-\pi}^{\pi}2(1-\cos t)F_s(t)^2\,\frac{dt}{2\pi}
 =\sum_{j\in\Z}(b_j-b_{j-1})^2=\frac2s.
\end{align*}
The cosine difference bound follows from
$\abs{\cos\theta-\cos(\theta-t)}\le2\abs{\sin(t/2)}$;
the integral identity is Parseval's identity applied to
$(1-e^{\mathrm{i}t})F_s(t)$. Exactly $2s$ successive differences have
modulus $1/s$. Since $0\le p_{s,N}\le2/C_s$ on the interval,
$0\le j_s\le c_s$. Its degree cap $D+2$ is even, so
\cref{lem:interval-sos} certifies $c_s-j_s$ with that cap.

Each $\omega$ is globally SOS, and $1-\omega=\varrho$ is interval
nonnegative with a degree-$D$ interval certificate. Telescoping
$1-\prod_{\ell\ne i}\omega(x_\ell)$, with SOS multipliers consisting of
preceding $\omega$ factors, gives an ordinary-module certificate. Hence
\[
 c_s-j_s(x_i)\prod_{\ell\ne i}\omega(x_\ell)
 =c_s(1-\prod_{\ell\ne i}\omega(x_\ell))
 +(c_s-j_s(x_i))\prod_{\ell\ne i}\omega(x_\ell)
\]
is certified through degree $v_bD+2\le r\le2r$. The product subtracted
from $c_s$ is SOS at the same degree, proving both inequalities.
\end{proof}

\begin{theorem}[Constrained ordinary-module rounding]\label{thm:constr-mod}
Assume \cref{ass:global-error}, $s\ge\max\{2,d_\infty\}$, even $N\ge2$,
and $r\ge wD+2d$. Every feasible family for
\eqref{eq:ext-mod-localizers} admits a probability law $\nu$ on $K$ with
\begin{equation}\label{eq:ext-mod-bound}
 \int f\,d\nu\le\sum_bL_b(f_b)+E+L_fH U^{\alpha/2}.
\end{equation}
Consequently $0\le\fstar-\rho_r^{\mathrm{c,mod}}\le E+L_fH U^{\alpha/2}$.
For fixed data and width, parameters can be chosen for all sufficiently
large $r$ so that
\[
 E=O(\log^3r/r^2),\qquad U=O(\log^2r/r^2),\qquad
 \fstar-\rho_r^{\mathrm{c,mod}}=O((\log r/r)^\alpha).
\]
\end{theorem}
\begin{proof}
Let $\Psi_b(x,u)=\prod_{i\in B_b}\Psi(x_i,u_i)$ and
$h_b(u)=L_b(\Psi_b(x,u))\ge0$. Since $\Psi_b$ is SOS of degree at most
$v_bD$, its conditional functional is positive on $(a+cg)^2$ and
nonnegative at $g$, for $g=g_{bj}$. Indeed the required degrees are at
most $v_bD+2d\le r$ and $v_bD+d\le2r$. The proof of
\eqref{eq:ext-conditional-violation} therefore applies to $\Psi_b$.
Multiply the certificate in \cref{lem:ext-difference} by the SOS
$\Psi_b$ and apply $L_b$ for each fixed $u$. Its degree is at most
$v_bD+2d\le r$. Integrating the resulting scalar inequality and using
\cref{lem:ext-displacement} gives
\begin{equation}\label{eq:ext-SOS-violation}
 \int(-g(u))_+^2h_b(u)\,d\arcs^{B_b}(u)
 \le\Abud(g)\sum_i A_i(g)
 L_b\Bigl(j_s(x_i)\prod_{\ell\ne i}\omega(x_\ell)\Bigr)
 \le c_s\Abud(g)^2.
\end{equation}
No measurable choice of SOS certificates is required: certificates
justify a scalar inequality at each fixed output before integration.

The source density has mass
$Z_b=L_b(\prod_i\omega(x_i))$ with $a_0^{v_b}\le Z_b\le1$.
For completeness, the upper bound follows from the telescope in
\cref{lem:ext-displacement}. For the lower bound telescope
$\prod_i\omega(x_i)-a_0^{v_b}$ into terms
$(\omega(x_i)-a_0)$ times preceding SOS mass factors and nonnegative
constant powers of $a_0$. Each $\omega-a_0\ge0$ has an interval
certificate through degree $D$. Both telescopes have degree at most
$v_bD$, within the available order.

By \cref{lem:residual-products}, the expansion of the signed density
$\bar h_b=L_b(\prod_i\bar\Psi(x_i,u_i))$ has zero-residual term $h_b$,
nonnegative singleton-residual terms, and higher terms whose sum is at
least $-\Delta_w$. Therefore its common correction satisfies
\[
 h_b^+=\frac{\bar h_b+\Delta_w}{1+\Delta_w}
 \ge\frac{h_b}{1+\Delta_w}.
\]
The nonnegative remainder
$e_b=h_b^+-h_b/(1+\Delta_w)$ has mass
$(1+\Delta_w-Z_b)/(1+\Delta_w)$. Since
$(-g)_+^2\le\chebnorm{g}^2$, \eqref{eq:ext-SOS-violation} gives
\[
 \int(-g)_+^2h_b^+\,d\arcs^{B_b}
 \le\frac{c_s\Abud(g)^2+\chebnorm{g}^2(1+\Delta_w-a_0^{v_b})}
 {1+\Delta_w}.
\]
The corrected bag laws
have exactly consistent separators and glue to a single box law
$\sigma$ by \cref{thm:mod}. That theorem gives objective error at most
$E$, and summing the last display gives $\int V\,d\sigma\le U$.
Apply \cref{lem:global-repair}, then take the hierarchy infimum.

To verify the rates, put $c_w=\max\{3,(w+1)/2\}$ and, for sufficiently
large $r$, choose
\[
 s=\left\lfloor\frac{r-2d}{2w(c_w+3)\log_2(r+2)}\right\rfloor,
 \qquad N=2\left\lceil\frac{c_w}{2}\log_2s\right\rceil.
\]
Then $s\ge\max\{2,d_\infty\}$, $N\ge2$ is even, and
$wD\le2ws(N+1)\le r-2d$. These are the parameters of
\cref{cor:mod-rate} with the additive reserve. They give
$\delta\le\tfrac12s^{-c_w}$,
$\eta=O_{d_\infty}(\log s/s^2)$, and
$\Delta_w=O_w((\log s)^{\max\{w-2,0\}}/s^3)=o_w(s^{-2})$
(the correction is zero for $w=1$). Thus
$1-a_0^{v_b}\le v_b\delta$, $U=O(s^{-2})$, and
$E=O(\log s/s^2)$. As $s=\Theta_w(r/\log r)$, the stated objective and
squared-violation estimates follow. For $0<\alpha\le1$ the objective
term is of smaller order than $(\log r/r)^\alpha$.
\end{proof}

\subsection{Certificates with positive slack}\label{sec:ext-certificates}

The constrained moment problems can fail strict feasibility, for example
when $K$ has empty interior. The full box blocks nevertheless imply equality
of the moment value and the certificate supremum.

Let $\mathcal C_r^{\mathrm{c,mod}}$ be the sum over bags of the cones of
polynomials
\[
 \sigma_{b0}+\sum_{i\in B_b}(1-x_i^2)\sigma_{bi}
       +\sum_jg_{bj}\sigma^g_{bj},
\]
where each $\sigma$ is SOS in $x_{B_b}$ and each summand has degree at most
$2r$.
Define $\mathcal C_r^{\mathrm{c,pre}}$ in the same way, using the summands
$q_I\sigma_{bI}$ and $g_{bj}q_I\sigma_{bjI}$ for all $I\subseteq B_b$.
All Gram bases are full; blocks with no allowed monomials are omitted.
These cones correspond exactly to
\eqref{eq:ext-mod-localizers} and \eqref{eq:ext-pre-localizers}.

\begin{proposition}[Constrained certificate supremum]\label{prop:constr-dual}
Suppose $K\ne\varnothing$, $r\ge1$, and $\deg f_b\le2r$ for every bag.
For $\mathcal C_r=\mathcal C_r^{\mathrm{c,mod}}$ or
$\mathcal C_r^{\mathrm{c,pre}}$, let $\rho_r$ be its corresponding moment
value. The moment minimum is attained, and
\[
 \sup\{\lambda:f-\lambda\in\mathcal C_r\}=\rho_r,
 \qquad f-\lambda\in\mathcal C_r\quad\text{for every }\lambda<\rho_r.
\]
No strict feasibility assumption is needed. The proposition does not assert
membership at $\lambda=\rho_r$.
\end{proposition}
\begin{proof}
Work in the finite-dimensional coefficient space
$W_r=\sum_b\Rle{x_{B_b}}{2r}$. The ordinary sparse box cone is contained
in either $\mathcal C_r$. Every bag monomial $x^\beta$ of degree at most
$2r$ splits as $x^\gamma x^\zeta$, with $\abs\gamma,\abs\zeta\le r$.
The identity
\[
 1\pm x^\beta=\tfrac12\bigl[(1-x^{2\gamma})+(1-x^{2\zeta})
                         +(x^\gamma\pm x^\zeta)^2\bigr]
\]
belongs to the ordinary box module at order $r$: each complement telescopes
into square multiples of individual box generators through degree $2r$.
Thus $1$ plus any coefficient perturbation whose absolute coefficients
sum to less than one is in $\mathcal C_r$, by taking nonnegative
combinations of $1$ and these identities. It follows that
$1\in\operatorname{int}_{W_r}\mathcal C_r$.

By \cref{lem:zero-decomposition}, normalized linear functionals on $W_r$
that are nonnegative on $\mathcal C_r$ correspond exactly to the local
moment families: separator equalities identify the different bag
representations of a polynomial. Evaluation at any point of $K$ makes this
set nonempty. It is closed and bounded by the box moment bounds in
\cref{lem:compact}, so its objective minimum $\rho_r$ is attained.

Fix $\lambda<\rho_r$ and suppose $f-\lambda\notin\mathcal C_r$.
Finite-dimensional separation from the interior of the convex cone gives
a nonzero linear functional $F$ with $F(q)\ge0$ for every
$q\in\mathcal C_r$ and $F(f-\lambda)\le0$. This separation does not require
closedness of the cone. Since $1$ is an interior point, $F(1)>0$.
Normalizing $F$ therefore gives a feasible moment functional with objective
at most $\lambda$, contradicting its strict inequality to $\rho_r$.
Hence every strict level is certified. Weak duality gives the reverse
inequality between the certificate supremum and $\rho_r$.
\end{proof}

In particular, let $E_r$ denote the last two terms of
\eqref{eq:ext-pre-bound}, or $E+L_fHU^{\alpha/2}$ in
\eqref{eq:ext-mod-bound}. The corresponding theorem and
\cref{prop:constr-dual} give
\[
 f-\fstar+E_r+\varepsilon\in\mathcal C_r
 \qquad\text{for every }\varepsilon>0.
\]
Thus the same rates apply to real certificates with arbitrarily small
positive slack, including problems with equality constraints or empty
interior. Boundary attainment remains a separate question.

The constants retain the sums of generator budgets, $H$, and $L_f$.
The exact-density proof introduces no additional tree disagreement term.
Feasible repair is an existence result and can involve a global nonconvex
projection problem. In the linear-error-bound class, the actual-measure
example of \cref{thm:affine-sharp} gives an inverse-order sparse gap even
with stronger local positivity; the power one is therefore sharp for this
class. It does not prove that the ordinary-module logarithm is necessary.

\begin{remark}[Prior methods and the geometric assumption]\label{rem:ext-prior}
Moment approximation followed by projection onto a feasible set is an
established method. Tran and Toh~\cite{tran2026-truncated-moment-sequences}
give dense preordering bounds with the same global Hölder exponent.
Heijmans-Kuryatnikova, Vera and
Zuluaga~\cite{heijmans2026-degree-bounds} transfer box certificates to
general semialgebraic sets by a polynomial lift. For a dense formulation
of the present global error-bound setting, after a fixed normalization of the generators,
combining their lift with the dense ordinary box rate gives a bound
$O((\log^{3/2}r/r)^\alpha)$ for fixed dimension and polynomial data; the direct displacement
estimate above improves this particular logarithmic composition.
This comparison concerns a dense certificate cone.
These comparisons do not establish priority for the refinement.
Korda, Magron and Ríos-Zertuche~\cite{korda2025-convergence-rates-sparsity}
use local geometric assumptions in their sparse general-domain theorem.
\Cref{ex:global-local} shows that local error bounds need not preserve the
global exponent, so these rates do not imply uniform improvement for all
such domains.
\end{remark}

\subsection{Finite-state labels}\label{sec:ext-finite}

Now let each bag have continuous variables $X_b$ and finite-state variables
$Z_b$, so $B_b=X_b\sqcup Z_b$. Every finite-state variable has an explicitly
listed finite domain, and running intersection holds for every variable.
Let $\mathcal A_b$ be the permitted assignments to $Z_b$, and assume
\[
 \mathcal A=\{a:a|_{Z_b}\in\mathcal A_b\text{ for every }b\}\ne\varnothing.
\]
All labels have the same continuous box. There are no additional continuous
constraints in this model. Its objective is
$f(a,x)=\sum_bf_{b,a|_{Z_b}}(x_{X_b})$, with polynomial local costs, and
$\fstar$ is its minimum over $\mathcal A\times[-1,1]^n$.
In this subsection $v_b=\abs{X_b}$ and $w=\max_b v_b$ are continuous sizes;
empty $X_b$ and $w=0$ are allowed. This explicitly overrides the standing
$w\ge1$, $r\ge1$ convention of \cref{sec:setting} when all variables are
finite-state. An empty product kernel is one, and $\arcs^\varnothing$
is the unit measure on the one-point empty product.

For each $a\in\mathcal A_b$, introduce
$L_{b,a}:\Rle{x_{X_b}}{2r}\to\R$, positive on either
$\Pre_r(X_b)$ or $\Mod_r(X_b)$, with
$\tau_{b,a}=L_{b,a}(1)\ge0$ and $\sum_a\tau_{b,a}=1$.
For an edge $e=bc$ put
$S_e^x=X_b\cap X_c$ and $S_e^z=Z_b\cap Z_c$. For every assignment $s$
to $S_e^z$ and every $p\in\Rle{x_{S_e^x}}{2r}$ require
\begin{equation}\label{eq:ext-label-separator}
 \sum_{a\in\mathcal A_b:\ a|_{S_e^z}=s}L_{b,a}(p)
 =\sum_{a\in\mathcal A_c:\ a|_{S_e^z}=s}L_{c,a}(p).
\end{equation}
An empty fiber has sum zero; an empty discrete separator has one empty
assignment. Equalities must include assignments whose fiber is absent on
one side. Equating individual full-bag labels would impose the wrong
condition. Denote the two moment values by
$\rho_r^{\mathrm{mix,pre}}$ and $\rho_r^{\mathrm{mix,mod}}$.

\begin{lemma}[Mass-scaled moment control]\label{lem:ext-mass}
If $L$ is ordinary-box-module-positive through degree $2r$ with mass
$\tau\ge0$, then $\abs{L(T_\beta)}\le\tau$ for $\abs\beta\le r$.
If $\tau=0$, the entire functional through degree $2r$ vanishes.
Both conclusions also hold under preordering positivity.
\end{lemma}
\begin{proof}
The proof of \cref{lem:cheb-moment} is homogeneous: a certificate for
$1-T_\beta^2$ through degree $2\abs\beta$ gives
$0\le L(T_\beta^2)\le\tau$, and the PSD form on $1,T_\beta$ gives
$\abs{L(T_\beta)}^2\le\tau L(T_\beta^2)\le\tau^2$.
If $\tau=0$, all these diagonal values are zero. The PSD moment form on
polynomials of degree at most $r$ has zero diagonal in its Chebyshev basis,
so every entry is zero by Cauchy--Schwarz. It follows that $L(pq)=0$ for
all $\deg p,\deg q\le r$. Every monomial through degree $2r$ is such a
product: split its exponent into two exponents of total degree at most
$r$. Thus all moments vanish. For an empty continuous bag the only moment
is the mass, so this conclusion also applies at order zero.
\end{proof}

Set
\[
 A_{\mathrm{mix}}=\sum_b\max_{a\in\mathcal A_b}\Abud(f_{b,a}),\quad
 C_{b,a}=\chebnorm{f_{b,a}-c_{b,a,0}},\quad
 C_{\mathrm{mix}}=\sum_b\max_{a\in\mathcal A_b}C_{b,a},
\]
where $c_{b,a,0}$ is the constant Chebyshev coefficient. Let $d$ and
$d_\infty$ bound the total and coordinate degrees of these costs, taking
both zero for constant costs. Set $\Gamma_0=\Delta_0=0$.

\begin{theorem}[Finite-state preordering rounding]\label{thm:finite-state}
For $w\ge1$ and $r\ge\max\{w,d\}$, put
$m=\lfloor r/w\rfloor+1$. Every feasible labelled preordering family has a
probability law $\nu$ supported on $\mathcal A\times[-1,1]^n$ such that
\[
 \Bigl\lvert\int f\,d\nu-\sum_{b,a}L_{b,a}(f_{b,a})\Bigr\rvert
 \le\frac3{\Dm}\sum_{b,a}\tau_{b,a}\Abud(f_{b,a})
 \le\frac{3A_{\mathrm{mix}}}{\Dm}.
\]
Hence $0\le\fstar-\rho_r^{\mathrm{mix,pre}}\le3A_{\mathrm{mix}}/\Dm$.
Empty continuous bags are permitted. If $w=0$, the hierarchy is exact at
order zero; this case uses neither $m$ nor a division by $w$.
\end{theorem}
\begin{proof}
For each label define
$h_{b,a}(u)=L_{b,a}(K_{m,X_b}(x,u))$. The kernel's degree is at most
$2v_b(m-1)\le2r$, and its preordering certificate gives $h_{b,a}\ge0$.
Normalization of the kernel gives $\int h_{b,a}\,d\arcs^{X_b}=\tau_{b,a}$.
After continuous marginalization and summation over the labels extending
$s$, the separator density is
\begin{equation}\label{eq:ext-pre-separator-density}
 \sum_{a:\ a|_{S_e^z}=s}L_{b,a}(K_{m,S_e^x}(x,u)).
\end{equation}
It agrees at $b,c$ by \eqref{eq:ext-label-separator}. Thus these are
consistent mixed bag probability laws. Apply \cref{lem:gluing} on the
products of finite spaces and intervals. Equivalently, use the explicit
conditional-density construction of \cref{rem:ker-gluing}, summing over
new finite labels and integrating over new continuous coordinates.
At zero separator density any conditional law can be chosen, since the
event is null. Every bag law assigns zero probability to forbidden local
labels; a finite union of these null events shows that the global law has
finite-state support in $\mathcal A$.

The Chebyshev multiplier calculation of \cref{thm:pre}, with
\cref{lem:ext-mass}, gives error at most
$3\tau_{b,a}\Abud(f_{b,a})/\Dm$ for each label. Sum and use bag mass
normalization. Constants depending on labels are preserved exactly. Empty
continuous bags have density $\tau_{b,a}$ and zero error. If $w=0$, all
functionals are scalar masses, \eqref{eq:ext-label-separator} gives
consistent finite bag probabilities already at order zero, and their
exact gluing preserves the whole objective. Feasible Dirac families give
weak duality and therefore exactness.
\end{proof}

\begin{theorem}[Finite-state ordinary-module rounding]\label{thm:finite-state-mod}
Suppose $w\ge1$, $s\ge\max\{2,d_\infty\}$, even $N\ge2$, and $r\ge wD$.
Every feasible labelled ordinary-module family admits a law $\nu$ supported
on $\mathcal A\times[-1,1]^n$ whose local labelled densities are
\begin{equation}\label{eq:ext-label-correction}
 h^+_{b,a}(u)=\frac{\bar h_{b,a}(u)+\Delta_w\tau_{b,a}}{1+\Delta_w},
 \qquad \bar h_{b,a}=L_{b,a}(\bar\Psi_{X_b}(x,u)),
\end{equation}
and whose objective satisfies
\begin{align}
 \Bigl\lvert\int f\,d\nu-\sum_{b,a}L_{b,a}(f_{b,a})\Bigr\rvert
 &\le\sum_{b,a}\tau_{b,a}
   \bigl[C_{b,a}\Gamma_{v_b}+\Delta_w\osc(f_{b,a})\bigr]
       \label{eq:ext-label-weighted}\\
 &\le C_{\mathrm{mix}}(\Gamma_w+2\Delta_w).\label{eq:ext-label-data}
\end{align}
The data-only bound \eqref{eq:ext-label-data} bounds
$0\le\fstar-\rho_r^{\mathrm{mix,mod}}$. For fixed $w,d_\infty$ it gives
$C_{\mathrm{mix}}O_{w,d_\infty}(\log^3r/r^2)$ using the parameters of
\cref{cor:mod-rate}. Zero-mass labels and empty continuous bags require no
normalization. If $w=0$, the ordinary labelled hierarchy is exact at order
zero, with the same conventions as \cref{thm:finite-state}.
\end{theorem}
\begin{proof}
The proof of \cref{lem:mod-signed} works with mass $\tau_{b,a}$: the
signed density has this mass, and marginalization deletes the removed
kernel factors exactly. Its coefficient error is at most
$\tau_{b,a}C_{b,a}\Gamma_{v_b}$ by \cref{lem:ext-mass}, since the
transformed and original polynomials have degree at most $v_bD\le r$.
Indeed $\deg f_{b,a}\le v_bd_\infty\le v_bs\le v_bD$ for nonempty bags;
empty bags have constant costs.

In the residual expansion, for a subset of size $j$ put
$G_J=\prod_{i\notin J}\Psi(x_i,u_i)$ and
$H_J=\prod_{i\in J}\varrho(x_i)$. The certificate and weighted
Cauchy--Schwarz proof of \cref{lem:residual-products} gives
\[
 \abs{L_{b,a}(G_JH_J)}\le\delta^jL_{b,a}(G_J)
 \le\tau_{b,a}\delta^j\Lambda^{v_b-j}.
\]
This proof never divides by the mass: the residual-square certificates
have degree at most $(v_b+j)D\le2r$, and the coefficient bound for $G_J$
uses degree $(v_b-j)D\le r$. The terms with $j=0,1$ are nonnegative.
Hence $\bar h_{b,a}\ge-\tau_{b,a}\Delta_{v_b}\ge-\tau_{b,a}\Delta_w$.
The correction \eqref{eq:ext-label-correction} is nonnegative and retains
mass $\tau_{b,a}$. A zero-mass functional vanishes by
\cref{lem:ext-mass}; an empty bag has $\bar h_{b,a}=h^+_{b,a}=\tau_{b,a}$.

For a separator label $s$, its corrected continuous density is exactly
\[
 \frac{\displaystyle
 \sum_{a:\ a|_{S_e^z}=s}L_{b,a}(\bar\Psi_{S_e^x}(x,u))
 +\Delta_w\sum_{a:\ a|_{S_e^z}=s}\tau_{b,a}}{1+\Delta_w}.
\]
Both sums agree at the endpoints by \eqref{eq:ext-label-separator}, the
second using $p=1$. Absent fibers remain zero. Thus the mixed laws glue
exactly as in \cref{thm:finite-state}, including null separators and
permitted support. Adding a bare constant to each label would generally
alter separator probabilities when the fibers have different sizes; the
mass factor in \eqref{eq:ext-label-correction} is essential.

Let $I^+_{b,a}=\int f_{b,a}h^+_{b,a}\,d\arcs^{X_b}$ and define
$\bar I_{b,a}$ with $\bar h_{b,a}$. Rearranging the correction gives
\[
 I^+_{b,a}-\bar I_{b,a}
 =\Delta_w\bigl(\tau_{b,a}\int f_{b,a}\,d\arcs^{X_b}-I^+_{b,a}\bigr).
\]
The two terms in parentheses integrate $f_{b,a}$ against nonnegative
measures of the same mass, so their difference is bounded by
$\tau_{b,a}\osc(f_{b,a})$. Combine this with the coefficient error and
sum to obtain \eqref{eq:ext-label-weighted}. Since
$\sum_a\tau_{b,a}=1$, $\osc(f_{b,a})\le2C_{b,a}$, and
$\Gamma_{v_b}\le\Gamma_w$, the data-only bound follows. It is uniform
over feasible families, so it can be used before taking the infimum.
The parameter proof of \cref{cor:mod-rate} applies unchanged. For $w=0$
there is no kernel parameter choice; the scalar-mass gluing proof gives
exactness at order zero directly.
\end{proof}

If $C_{\mathrm{mix}}=0$, all local costs depend only on their finite labels,
and the hierarchy is exact whenever their constant moments are defined.
There is no extra label-count factor in these errors; explicit label lists
and their SDP blocks can nevertheless be large. The results require a
common continuous box for every label. They do not cover label-dependent
boxes or additional label-dependent continuous constraints.

\begin{proposition}[Real certificates after label pruning]\label{prop:finite-dual}
Delete every local label that has no extension in $\mathcal A$. This
preserves the labelled primal values and feasible objectives. For the
remaining finite SDP, at any order containing the objective, both labelled
box hierarchies have an attained primal value and an attained dual of the
same value. The certificate is an identity of labelled local polynomials
on the permitted assignment space.
\end{proposition}
\begin{proof}
The scalar masses of any feasible family satisfy consistent discrete bag
probabilities, so finite tree gluing gives a law on $\mathcal A$. A label
without a global extension therefore has zero mass. Its whole functional
vanishes by \cref{lem:ext-mass}; deleting its blocks and zero terms in the
equalities preserves feasibility and objectives, in both directions.

Choose a distribution positive on every element of the finite nonempty
set $\mathcal A$, independently of continuous product uniform measure.
Every retained label has positive mass. A nonzero local polynomial square,
weighted by any retained box-generator product, has strictly positive
integral on the open box. Thus every retained labelled PSD block is
positive definite; scalar blocks of empty continuous bags are positive
as well. All separator and normalization equalities hold. This is primal
Slater in the pruned finite SDP. Feasible moments are bounded by their
masses: for monomials of degree at most $r$, a telescope for
$1-x^{2\beta}$ gives $0\le L_{b,a}(x^{2\beta})\le\tau_{b,a}$; splitting
any exponent through degree $2r$ and applying Cauchy--Schwarz bounds its
moment by $\tau_{b,a}\le1$. Consequently the feasible set is closed and
bounded, so the primal optimum is finite and attained. Finite SDP Slater
duality, as used in \cref{lem:slater}, gives dual attainment and equality.

To describe the identity, orient every edge and write
$\epsilon_{b,e}=+1$ at its first endpoint, $-1$ at its second. Dual
multipliers are scalars $\theta_b$ for normalization and separator
polynomials $\phi_{e,s}(x_{S_e^x})$ of degree at most $2r$ for each
separator label $s$. Write $\rho_r^{\mathrm{mix}}$ for the value of the
chosen labelled hierarchy. At the optimum they satisfy, for every retained label,
\[
 f_{b,a}-\theta_b-
 \sum_{e\ni b}\epsilon_{b,e}\phi_{e,a|_{S_e^z}}
 \in\Pre_r(X_b)\quad\text{or}\quad\Mod_r(X_b),
 \qquad \sum_b\theta_b=\rho_r^{\mathrm{mix}}.
\]
For each global permitted assignment the separator terms cancel when
these identities are summed. This is the asserted labelled certificate;
it is not an unlabelled polynomial identity in surrogate integer
variables. No attainment for the unpruned SDP is needed.
\end{proof}

Finite-state tree gluing and labelled formulations are established tools;
our use of them is a quantitative consequence of the kernels and the
mass-scaled common correction. Pruning justifies the real finite-SDP
certificate statement. It supplies no rational bit-size bound for the
labelled cone, and neither labelled theorem is a general MINLP runtime
claim.
